% ID: SEL-D0038-OA-E04
% ESTADO: cotejada. Segunda pasada visual y borrador compilado revisados.
% Fuente: Selectividad/matematicas_aplicadas/Ua3MATEMÁTICAS APLICADAS examen.pdf
% SHA-256: 4440c15a31a77ee9e052a0668923508e117795b369d569bc02d24b7937366f7a
% Página: 1; opción A; ejercicio original 4.

\begin{enumerate}[label=\alph*)]
\item \textbf{(1.5 puntos)} En una muestra aleatoria de 100 botellas de agua mineral se encontró un contenido medio de 48 cl. Sabiendo que la variable ``contenido de agua en una botella'' sigue una ley Normal con desviación típica 5 cl, determine un intervalo de confianza para la media poblacional, con un nivel de confianza del 95\%.
\item \textbf{(1 punto)} \textquestiondown{}Qué tamaño muestral mínimo debería considerarse para estimar esta media con el mismo nivel de confianza y un error inferior a 0.5 cl?
\end{enumerate}
